\begin{abstract}
\normalsize

针对圆柱形药材热风干燥中温湿分布不均、表面及平均状态难以代表全域干燥程度的问题，
本文建立一维径向热湿传递模型。在常热物性模型基础上，进一步考虑
状态依赖物性、全域达标判据及实测径向收缩，采用
\textbf{环形有限体积法}进行空间离散，并以
\textbf{自适应隐式BDF方法}推进时间，计算规定时刻与位置的
温度和含水率，分析内部水分迁移与尺寸变化对烘干进程的影响。

\textbf{针对问题一，}在常热物性条件下分别建立径向热传导
与非线性水分扩散模型，以时变环境驱动表面交换，求得预热
阶段的温湿分布。预热 \(1800~\mathrm{s}\) 后，中心温度升至
\(33.58~^\circ\mathrm{C}\)，含水率仍接近初值，而含水率
下降超过 \(5\%\) 的区域仅从表面向内延伸约
\(5.69~\mathrm{mm}\)。按初始物性估算，热扩散与水分扩散
特征时间分别约为 \(39.5~\mathrm{min}\) 和
\(22.5~\mathrm{h}\)，后者约为前者的 \(34.2\) 倍。
这与预热时温度已传入中心、明显失水仍集中于外层的
计算结果一致。

\textbf{针对问题二，}依据含水率更新体积热容与导热系数，
并结合温度、含水率更新水分扩散系数，从题定初态实现变物性热湿联立
求解。\(3~\mathrm{h}\) 时，中心与表面温差仅为
\(0.12~^\circ\mathrm{C}\)，含水率差仍达
\(0.76~\mathrm{kg/kg}\)，表明温度趋于均匀时内部仍存在
明显水分梯度。进一步设置初态物性冻结对照，发现相较于
变物性组，冻结组预测的中心含水率更高、表面含水率更低，
该时刻的体积平均含水率高 \textbf{8.39\%}。同幅度参数扰动下，
前三小时的平均含水率对表面传质系数更敏感，最湿点含水率
则对内部扩散系数更敏感。

\textbf{针对问题三，}考虑环境记录仅覆盖前 \(4~\mathrm{h}\)
的情况，以观测末 \(60~\mathrm{min}\) 的环境均值延拓后续
边界，采用\textbf{全域阈值事件定位}，求取最大含水率下降
穿越 \(0.15~\mathrm{kg/kg}\) 的临界时刻，并继续积分核验
严格达标状态。计算得到固定尺寸模型的烘干临界时间为
\textbf{57.4796~h}。体积平均含水率在 \(36.03~\mathrm{h}\)
达到阈值时，仍有 \textbf{56.66\%} 的体积未达标，若据此
终止将提前约 \(21.45~\mathrm{h}\)，说明平均值不能代替
全域达标判据。最湿位置完成最后 \(5\%\) 目标含水率降幅
仍需 \(37.20~\mathrm{h}\)，占总时长的 \textbf{64.71\%}；
扩散系数与表面传质系数的同幅度扰动对照进一步表明，
在所考察范围内，末段耗时对扩散能力的响应明显更强，
支持干燥末段主要受内部扩散限制的判断。

\textbf{针对问题四，}结合实测半径与附录4物性，通过
\textbf{材料坐标变换}将径向收缩区域映射为固定计算区间，
并设置\textbf{同物性、同环境的固定半径对照}，考察给定
物性条件下的几何收缩作用。在相同长期环境方案下，
收缩模型的临界时间为 \textbf{51.0934~h}，固定半径对照
为 \(129.8622~\mathrm{h}\)，引入实测收缩后缩短
\textbf{60.66\%}。
进一步比较半径与含水率历程发现，\(24~\mathrm{h}\) 时
尺寸已基本稳定，但内部仍需 \(27.09~\mathrm{h}\) 才能
全域达标，说明外形稳定不能直接作为烘干结束依据。

通过温度解析参照、制造解、独立方法对照与网格加密检验数值求解，
通过事件前后状态核验临界时刻，并以环境、传质参数与收缩输入扰动
考察预测结果的敏感性。所得临界时间为给定物性、环境与几何假设
下的径向模型预测，可为内部湿芯识别和烘干终止判断提供参考。

\vspace{1em}
\noindent\textbf{关键词：}
热湿耦合；有限体积法；全域达标判据；事件定位；径向收缩

\end{abstract}
